\section{RNN 隐藏状态的可解释性}

RNN 的一个有趣特性是，我们可以直接观察隐藏状态中各个神经元的激活值，并发现其中一些具有高度可解释的语义。

\begin{figure}[H]
\centering
\includegraphics[width=0.95\textwidth]{slides-images/slide-028.jpg}
\caption{RNN 隐藏状态可视化：使用 $\tanh$ 激活函数，值域为 $[-1, 1]$，红色表示接近 $-1$，蓝色表示接近 $+1$\protect\footnotemark}
\end{figure}
\footnotetext{Slides 第29页。}

讲者展示了在代码和文本数据上训练的 RNN 中发现的几种可解释的隐藏单元：

\begin{figure}[H]
\centering
\includegraphics[width=0.95\textwidth]{slides-images/slide-029.jpg}
\caption{引号检测器：该神经元在引号开始时激活，引号结束时关闭，准确追踪引号的开闭状态\protect\footnotemark}
\end{figure}
\footnotetext{Slides 第30页。}

\begin{figure}[H]
\centering
\includegraphics[width=0.95\textwidth]{slides-images/slide-030.jpg}
\caption{行长度追踪器：该神经元的激活值随行的进展逐渐变化，在接近换行符时达到极值\protect\footnotemark}
\end{figure}
\footnotetext{Slides 第31页。}

\begin{figure}[H]
\centering
\includegraphics[width=0.95\textwidth]{slides-images/slide-031.jpg}
\caption{if 语句检测器：在代码中，该神经元在进入 if 语句块时激活\protect\footnotemark}
\end{figure}
\footnotetext{Slides 第32页。}

\begin{figure}[H]
\centering
\includegraphics[width=0.95\textwidth]{slides-images/slide-033.jpg}
\caption{代码缩进深度追踪器：随着代码嵌套层次加深，该神经元的激活值逐步增大\protect\footnotemark}
\end{figure}
\footnotetext{Slides 第34页。}

\begin{knowledgebox}{RNN 的自发特征学习}
这些可解释的隐藏单元\textbf{从未被明确设计或监督}——它们完全是模型在训练过程中自发学会的。这与我们手工构造的玩具 RNN 形成了有趣的对比：
\begin{itemize}
    \item 玩具示例中，我们手动指定隐藏状态追踪``当前值''和``前一步值''
    \item 训练好的 RNN 自动学会追踪引号状态、代码深度等更复杂的概念
    \item 但并非所有隐藏单元都可解释——大部分单元的激活模式看起来是随机的
\end{itemize}
\end{knowledgebox}

\subsection{本章小结}

RNN 隐藏状态中的某些神经元能够自发地学习到有意义的特征，如引号追踪、行长度计数、代码嵌套深度等。这种可解释性是 RNN 架构的一个独特优势。

%% ========== RNN 优缺点 ==========

